\chapter{Experimental Evaluation}
\label{ch:experiments}

This chapter presents the empirical evaluation of HSP-Agile against three baseline
methods and three ablated variants.  We structure the evaluation around four research
questions (Section~\ref{sec:exp:rq}), describe the experimental setup and benchmark
(Sections~\ref{sec:exp:setup}--\ref{sec:exp:benchmark}), define the compared methods
and metrics (Sections~\ref{sec:exp:methods}--\ref{sec:exp:metrics}), and then report
results for each research question in turn
(Sections~\ref{sec:exp:rq1}--\ref{sec:exp:rq4}).  We close with a sensitivity
analysis (Section~\ref{sec:exp:sensitivity}) and a structured threats-to-validity
discussion (Section~\ref{sec:exp:threats}).

%=============================================================================
\section{Research Questions}
\label{sec:exp:rq}
%=============================================================================

The evaluation is designed to answer the following four research questions.

\paragraph{RQ1 --- Strict Correctness and Formal Conformance.}
Does HSP-Agile achieve higher strict formal conformance than the one-shot (B1) and
test-feedback (B2) baselines?  The formal conformance metric $\mathit{Conf}(c,t)$
(Definition~\ref{def:conformance}) captures the fraction of Z3-generated evaluation
cases that a candidate implementation satisfies; it is a more granular quality signal
than the binary pass/fail metric used by B1 and B2.  We are particularly interested
in whether the CEGIS repair loop and Z3-based case generation together push
conformance beyond what iterative unit-test feedback achieves.

\paragraph{RQ2 --- Operational Efficiency.}
What is the latency cost of HSP-Agile relative to simpler baselines?  Each HSP-Agile
run invokes Z3 multiple times, queries the LLM in a repair loop, and evaluates pattern
guards, all of which add wall-clock time.  We measure mean per-task latency for all
seven modes and quantify the overhead of each pipeline component.  The goal is to
determine whether the quality gains from RQ1 justify the additional cost in a
realistic sprint-cycle deployment scenario.

\paragraph{RQ3 --- Component Contribution (Ablation).}
How much does each pipeline component individually contribute to formal conformance?
We isolate the formal checker (A1), the pattern guard (A2), and the repair loop (A3)
by disabling each component in turn, holding all other parameters fixed.  The ablation
decomposition lets us rank the contributions and understand the interaction effects
between components.

\paragraph{RQ4 --- Defect Prevention Capability.}
Does HSP-Agile detect more faulty candidates (higher Positive Detection Rate, PDR)
with fewer false accepts (lower False Accept Rate, FAR) than the baselines?  PDR and
FAR measure the pipeline's utility as a quality gate: a system deployed in a CI
pipeline must catch real faults while not blocking correct submissions~\citep{demillo1978mutation,offutt1992mutation}.
Because the full prevention evaluation requires generating and classifying large mutant
populations, we report preliminary results for three modes (B1, B2, M) and
characterise the remaining evaluation as ongoing work.

%=============================================================================
\section{Experimental Setup}
\label{sec:exp:setup}
%=============================================================================

\subsection{Hardware and Software Environment}

All experiments were executed on a desktop workstation running Windows 10
(build 26220) with a 6-core processor.  Tasks were dispatched in parallel across
all six cores using Python's \texttt{multiprocessing} module, with one worker per
core.  The Python runtime was version 3.12.  The Z3 SMT solver version 4.16~\citep{demoura2008z3,barrett2018smt}
was invoked via the \texttt{z3-solver} Python binding; all Z3 calls used a per-call
timeout of 10 seconds, after which the solver returned \texttt{unknown} and the
pipeline treated the call as a non-constraining outcome.

\subsection{Language Model Configuration}

All LLM-based modes (B1, B2, M, A1, A2, A3) used Qwen2.5-27B-Instruct
(\textsf{Qwen-27B}) accessed through an OpenAI-compatible LLM gateway.  Temperature was set to 0.7 for all
generation calls and to 0.0 for all repair calls, following the prompting strategy
described in Chapter~\ref{ch:method}.  The maximum token budget per LLM call was
4096 tokens.  The budget parameter $K = 3$ was fixed for all LLM-based modes: in
mode M, the repair loop may invoke the LLM up to three times per task, consuming
between 1 and 3 API calls depending on when the loop terminates.

\subsection{Experiment Identifier and Scale}

The experiment was conducted under the run identifier
\texttt{run\_hard\_full\_parallel\_v1}.  A total of 840 runs were executed: 7 modes
$\times$ 120 tasks, each run processed independently by a dedicated worker process.
Run artefacts (prompts, LLM responses, Z3 query logs, evaluation traces, timing
records) were serialised to disk in JSONL format for reproducibility~\citep{ammann2016testing}.
All random-number generation used fixed seed 42 for task generation and prompt
construction; the LLM API itself is not seeded, but temperature 0.0 for repair
calls makes those responses near-deterministic in practice.

\subsection{Reproducibility Measures}

Determinism was enforced at three levels.  First, task generation used a fixed
random seed, ensuring the same 120 tasks for every replication.  Second, the
evaluation harness (Z3 case generation, mutation testing) is fully deterministic
given the same task and candidate.  Third, statistical comparisons use
non-parametric tests (Wilcoxon signed-rank and Mann-Whitney U) that do not assume
normality, making them robust to residual LLM non-determinism in one-shot calls.

%=============================================================================
\section{Benchmark Construction}
\label{sec:exp:benchmark}
%=============================================================================

\subsection{Design Rationale}

A central challenge in evaluating specification-conformance tools is that
off-the-shelf software repositories do not contain Agile-SOFL\slash FSF
specifications; such artefacts must be constructed.  We required a benchmark with
three properties: (i)~\emph{discriminative power}---tasks should be hard enough
that baseline methods fail non-trivially, otherwise ceiling effects mask quality
differences; (ii)~\emph{formal validity}---every task must be satisfiable by at
least one correct implementation, verifiable by Z3; and
(iii)~\emph{diversity}---tasks should cover varied input combinations, scenario
orderings, and boundary conditions to prevent overfitting by any single strategy.
Existing formal-synthesis benchmarks such as DafnyBench~\citep{loughman2024dafnybench}
target annotation-heavy verification languages rather than FSF scenario semantics,
motivating a domain-specific corpus.

\subsection{Base Tasks}

We began with 20 base Agile-SOFL tasks extracted by manual inspection from published
example specifications~\citep{liu2014sofl,li2022agilesofl}.  Each base task defines a small module with 3--5 inputs, 2--3
outputs, and 5--8 FSF scenarios that describe the required input-output relationship
in each region of the input space.

\subsection{Hard Synthetic Task Generation}

The 20 base tasks alone are insufficient for a statistically powered comparison.  We
therefore developed the \emph{HardTaskGenerator}, a programmatic generator that
produces synthetic FSF tasks with controlled difficulty.  Each generated task has
exactly 5 integer inputs $(a, b, c, d, e) \in [0, 100]^5$, 3 outputs
(\textit{risk}, \textit{action}, \textit{score}), and exactly 8 ordered FSF
scenarios.  Difficulty is modulated by two design choices.

\emph{Precedence-sensitive guard regions.}  Each scenario's pre-condition guard is
constructed so that it overlaps with at least one other guard in the input space.
Correct scenario evaluation therefore requires strict priority ordering: an
implementation that evaluates guards in the wrong order will produce correct outputs
on most inputs but fail on boundary inputs that fall in the overlap region.  This
design directly targets the most common fault pattern (scenario-order
violation, SOV) identified in the fault taxonomy of
Chapter~\ref{ch:formalization}.

\emph{Z3 satisfiability filter.}  Before a generated task is admitted to the
benchmark, Z3 verifies that all 8 scenarios are individually satisfiable: for each
scenario $s_i$, the conjunction of its guard and output constraints must have at
least one model.  Tasks failing this check are discarded.  This prevents benchmark
tasks that are logically vacuous and ensures every benchmark task has a ground-truth
correct implementation.

Figure~\ref{fig:benchmark-pipeline} (below) illustrates the full construction
pipeline.

\begin{figure}[t]
  \centering
  \tikzinput{diagrams/tikz/benchmark_pipeline}
  \caption{Benchmark construction pipeline.  Hard synthetic tasks are generated with
    precedence-sensitive FSF scenarios, validated by Z3 for scenario satisfiability,
    and assembled into the final 120-task benchmark used across all seven evaluation
    modes (840 total runs).}
  \label{fig:benchmark-pipeline}
\end{figure}

The generator produced 180 candidate tasks, of which 60 were discarded by the Z3
satisfiability filter, leaving the final benchmark of 120 hard synthetic tasks.  The
base tasks were retained as a development set but were not included in the primary
evaluation to avoid potential information leakage between the prompt-engineering
phase and the evaluation phase.

\subsection{Illustrative Task Example}

To ground the benchmark description, the following illustrative task models a financial risk classifier with
inputs \texttt{a} (credit score, $[0,100]$), \texttt{b} (debt ratio, $[0,100]$),
\texttt{c} (income volatility, $[0,100]$), \texttt{d} (asset coverage, $[0,100]$),
and \texttt{e} (market exposure, $[0,100]$).  Outputs are \texttt{risk} $\in
\{0,1,2,3\}$ (risk tier), \texttt{action} $\in \{0,1,2\}$ (recommended action),
and \texttt{score} $\in [0,100]$ (aggregate score).  The eight FSF scenarios are
ordered from most-specific to least-specific:

\begin{enumerate}
  \item High credit AND low debt AND low volatility $\Rightarrow$ (risk=0, action=0, score=90)
  \item High credit AND low debt $\Rightarrow$ (risk=0, action=0, score=75)
  \item High credit AND moderate debt $\Rightarrow$ (risk=1, action=1, score=60)
  \item Low credit AND high debt AND high volatility $\Rightarrow$ (risk=3, action=2, score=15)
  \item Low credit AND high debt $\Rightarrow$ (risk=2, action=2, score=30)
  \item Low credit $\Rightarrow$ (risk=2, action=1, score=40)
  \item Moderate credit AND high exposure $\Rightarrow$ (risk=1, action=1, score=55)
  \item (default) $\Rightarrow$ (risk=1, action=0, score=65)
\end{enumerate}

Scenario 1 and Scenario 2 have overlapping guards (any input satisfying Scenario 1
also satisfies Scenario 2's guard).  An implementation that evaluates guards in
reverse order (Scenario 2 before Scenario 1) will produce \texttt{score}=75 instead
of 90 on inputs in the Scenario 1 region---a scenario-order violation that
the pattern guard is designed to detect.  Z3 can generate inputs precisely in this
overlap region (e.g., \texttt{a=85, b=15, c=10, d=80, e=20}), making such violations
detectable during the repair loop.

This example illustrates why the hard benchmark is discriminating: a one-shot LLM
generation may produce code that handles 7 of the 8 scenarios correctly but fails on
the Scenario 1/2 overlap because the LLM did not fully parse the priority ordering
from the specification text~\citep{chen2021codex,dakhel2023copilot}.  The CEGIS repair loop, armed with a Z3-generated
boundary input, can correct this in a single repair iteration~\citep{solarlezama2008sketching,lynn2024verifierloop}.

\subsection{Benchmark Characteristics}

Table~\ref{tab:benchmark-stats} summarises the properties of the 120-task benchmark;
the mean guard region overlap rate is 1.8 scenarios per input point, confirming that
correct priority evaluation is non-trivial on most inputs.

\begin{table}[t]
  \caption{Summary statistics of the 120-task hard benchmark.}
  \label{tab:benchmark-stats}
  \centering
  \begin{tabular}{lc}
    \toprule
    Property & Value \\
    \midrule
    Total tasks & 120 \\
    Inputs per task & 5 \\
    Outputs per task & 3 \\
    Scenarios per task & 8 \\
    Mean guard-overlap rate (scenarios/point) & 1.8 \\
    Tasks rejected by Z3 filter & 60 \\
    Tasks accepted & 120 \\
    Evaluation cases per task (strict\_eval) & 64 \\
    \bottomrule
  \end{tabular}
\end{table}

%=============================================================================
\section{Compared Methods}
\label{sec:exp:methods}
%=============================================================================

We compare seven evaluation modes.  The modes span the design space from a pure
formal reference implementation (B0) through increasingly capable LLM-based methods
(B1, B2) to the full HSP-Agile pipeline (M) and three ablated variants (A1, A2, A3)
that systematically remove one component each.

\begin{table}[t]
  \caption{Compared methods and their component configurations.
           Formal = formal check; Guard = pattern guard; Repair = repair loop.}
  \label{tab:modes}
  \centering
  \thesistable{%
    \footnotesize
    \renewcommand{\arraystretch}{1.15}
    \begin{tabular*}{\linewidth}{@{\extracolsep{\fill}}l l c c c c@{}}
      \toprule
      Mode & Description & Formal & Guard & Repair \\
      \midrule
      B0 & Reference implementation   & Yes$^\dagger$ & No  & No           \\
      B1 & One-shot LLM               & No            & No  & No           \\
      B2 & LLM + test-feedback        & No            & No  & Yes (tests)  \\
      \textbf{M}  & \textbf{Full HSP-Agile} & \textbf{Yes (Z3)} & \textbf{Yes} & \textbf{Yes (CEGIS)} \\
      A1 & Ablation: no formal check  & No            & Yes & Yes          \\
      A2 & Ablation: no pattern guard & Yes (Z3)      & No  & Yes (CEGIS)  \\
      A3 & Ablation: no repair loop   & Yes (Z3)      & Yes & No           \\
      \bottomrule
    \end{tabular*}%
  }
  \par\smallskip
  \raggedright\footnotesize $^\dagger$B0 uses formal checking \emph{for evaluation only};
  the reference implementation is pre-written from the task specification, not
  generated by a solver.
\end{table}

\paragraph{B0 --- Reference Implementation.}
B0 provides a formal upper-bound reference.  For each task, the reference
implementation is generated directly from the FSF specification by a deterministic
template renderer that encodes each scenario as a guarded branch in priority order.
B0 involves no LLM and no iterative loop; it executes in under 10 milliseconds.
The formal checker is applied \emph{ex post} for evaluation purposes to compute
$\mathit{Conf}(\cdot)$, but was not used to produce the implementation.  B0 answers
the question: ``How well does a formally correct but algorithmically na\"{\i}ve
template implementation perform on hard tasks?''

\paragraph{B1 --- One-Shot LLM.}
B1 makes a single LLM API call with a carefully constructed prompt containing the
full FSF specification.  The response is parsed, syntax-checked, and passed directly
to the evaluator.  No feedback loop is applied.  B1 represents the minimal
LLM-based approach and establishes how much capability a single prompt provides.

\paragraph{B2 --- LLM with Test-Feedback Loop.}
B2 extends B1 with an iterative test-feedback loop analogous to standard
test-driven code generation~\citep{zhu1997specification}.  After each LLM generation, a set of
unit tests derived from the FSF pre/post-conditions is executed; if any test fails,
the failure message and the current candidate are appended to the prompt and a new
LLM call is made.  The loop terminates when all tests pass or the budget $K=3$ is
exhausted.  B2 uses no Z3 and no pattern guard; its feedback signal is purely from
unit-test outcomes.

\paragraph{M --- Full HSP-Agile Pipeline.}
Mode M runs the complete pipeline as described in
Chapters~\ref{ch:method} and~\ref{ch:implementation}: the LLM generates an initial
candidate; Z3 generates counterexample inputs; the pattern guard screens for
scenario-order violations and guard-overlap faults; if the candidate fails either
check, a CEGIS repair loop invokes the LLM again with the counterexample as
additional context.  The loop repeats until the candidate passes all checks or $K=3$
is reached, at which point the best-scoring candidate is returned.

\paragraph{A1 --- Ablation: No Formal Check.}
A1 disables Z3-based case generation.  The formal conformance checker is replaced
by the same unit-test harness used in B2, so the feedback signal comes from
unit tests rather than SMT-generated counterexamples.  The pattern guard and repair
loop remain active.  This isolates the contribution of Z3-guided feedback.

\paragraph{A2 --- Ablation: No Pattern Guard.}
A2 disables the pattern guard.  Candidates are evaluated by the formal checker
(Z3) and passed directly to the repair loop if they fail; no structural pattern
screening is applied.  This isolates the contribution of the pattern guard.

\paragraph{A3 --- Ablation: No Repair Loop.}
A3 disables the repair loop entirely.  The LLM is called exactly once (as in B1),
but the Z3-based formal checker and pattern guard \emph{are} applied to evaluate the
resulting candidate.  The evaluation verdict is computed but no corrective feedback
is sent.  This isolates the contribution of the CEGIS repair loop.

%=============================================================================
\section{Evaluation Metrics and Statistical Protocol}
\label{sec:exp:metrics}
%=============================================================================

\subsection{Primary Metrics}

\paragraph{Strict Success Rate.}
For a task $t$ and mode $m$, the run is a \emph{strict success} if the returned
candidate implementation passes every one of the 64 Z3-generated evaluation cases
(i.e., the strict evaluation set used for final scoring).  The strict success rate
is the fraction of the 120 tasks for which this holds.  This is a binary, all-or-nothing
metric; it is the most demanding correctness criterion.

\paragraph{Strict Formal Conformance.}
The formal conformance $\mathit{Conf}(c, t) \in [0,1]$ (Definition~\ref{def:conformance})
is the fraction of the 64 strict evaluation cases that the candidate $c$ satisfies on
task $t$.  The reported metric is the mean $\overline{\mathit{Conf}}$ across all 120
tasks, providing a continuous quality signal that discriminates between methods even
when strict success rates are similar.

\paragraph{Mutation Kill Rate.}
Given a candidate $c$ for task $t$, we inject $n_\mu$ syntactic mutants (scenario
reordering, guard negation, output perturbation) and check whether the formal
checker correctly classifies each mutant as non-conforming.  The mutation kill rate
$\mu_k$ is the fraction of mutants killed; it measures the \emph{precision} of the
formal checker rather than the quality of the generated code.  Values near 100\%
indicate the checker is a reliable fault oracle; lower values indicate blind spots.

\paragraph{LLM Call Count.}
The mean number of LLM API calls made per task.  For B1 this is always 1.0; for
M it ranges from 1 to 3 depending on when the repair loop terminates, with a mean
of 2.08.

\paragraph{Latency.}
Wall-clock time from task start to final evaluation verdict, measured in milliseconds
and averaged over the 120 tasks per mode.  Latency includes LLM API round-trip time,
Z3 solving time, pattern guard evaluation time, and Python orchestration overhead.

\paragraph{PDR and FAR (Prevention Metrics).}
The Positive Detection Rate (PDR) is the fraction of faulty candidates (mutants) that
the pipeline correctly rejects before they would be accepted.  The False Accept Rate
(FAR) is the fraction of correct candidates incorrectly rejected.  Both metrics are
defined formally in Section~\ref{sec:form:metrics}; for the prevention evaluation in
Section~\ref{sec:exp:rq4} we report only PDR due to the preliminary nature of the
available data.

\subsection{Statistical Protocol}

All pairwise comparisons involving per-task binary outcomes (strict success) use the
\emph{Wilcoxon signed-rank test}, which is appropriate for paired non-normal
distributions.  Comparisons of latency distributions, which are unpaired and
right-skewed, use the \emph{Mann-Whitney U test}.  Effect sizes are reported as
rank-biserial correlation $r$ where applicable; values $|r| < 0.1$ are interpreted
as negligible, $0.1 \le |r| < 0.3$ as small, $0.3 \le |r| < 0.5$ as medium, and
$|r| \ge 0.5$ as large.  The significance threshold is $\alpha = 0.05$ with no
multiple-comparison correction applied within each research question (comparisons are
few and pre-specified).

%=============================================================================
\section{RQ1: Strict Correctness and Formal Conformance}
\label{sec:exp:rq1}
%=============================================================================

\subsection{Main Results}

Table~\ref{tab:main-results} presents the primary experimental results for all seven
evaluation modes across 120 tasks each.  Figures~\ref{fig:conf-distribution},
\ref{fig:conf-cdf}, \ref{fig:summary-metrics}, \ref{fig:llm-calls},
and~\ref{fig:radar-metrics} visualise the per-task conformance distributions,
aggregate metric profile, LLM-call cost, and multi-metric trade-offs.

\begin{table}[t]
  \caption{Main comparison results ($n = 120$ tasks per mode). Best formal conformance
    shown in bold. Strict = strict success rate; Conf.\ = formal conformance;
    Kill = mutation kill rate; Fail.\ = mean strict evaluation failures per task
    (out of 64).}
  \label{tab:main-results}
  \centering
  \thesistable{%
    \footnotesize
    \renewcommand{\arraystretch}{1.12}
    \begin{tabular*}{\linewidth}{@{\extracolsep{\fill}}l r r r r r@{}}
      \toprule
      Mode & Strict & Conf. & Kill & Lat.\ (ms) & Fail. \\
      \midrule
      B0   & 27.5\%          & 27.5\%        & 79.0\%       & 7.5           & 0.00 \\
      B1   & 26.7\%          & 84.2\%        & 62.3\%       & 9{,}564       & 1.76 \\
      B2   & 26.7\%          & 88.0\%        & 61.4\%       & 10{,}769      & 1.85 \\
      \textbf{M}  & 25.0\%   & \textbf{90.4\%} & 60.6\%     & 43{,}659      & 1.98 \\
      A1   & 26.7\%          & 86.5\%        & 61.9\%       & 41{,}689      & 1.80 \\
      A2   & 26.7\%          & \textbf{90.9\%}& 60.1\%     & 12{,}080      & 1.97 \\
      A3   & 26.7\%          & 84.1\%        & 62.7\%       & 9{,}314       & 1.77 \\
      \bottomrule
    \end{tabular*}%
  }
\end{table}

\begin{figure}[t]
  \centering
  \widefig{figures/mode_distribution_strict_conformance.pdf}
  \caption{Distribution of per-task strict formal conformance ($\mathit{Conf}(c,t)$)
    across all 120 tasks for each evaluation mode.  Mode M achieves the rightmost
    distribution mass among the fully-specified modes, while B0's conformance is
    concentrated at lower values due to the hard benchmark's precedence-sensitive
    design.}
  \label{fig:conf-distribution}
\end{figure}

\begin{figure}[t]
  \centering
  \widefig{figures/formal_conformance_cdf.pdf}
  \caption{Empirical cumulative distribution of strict formal conformance by mode.
    Mode~M and A2 dominate the upper tail; B0's CDF rises steeply at low conformance
    values, reflecting template-renderer failures on precedence-sensitive tasks.}
  \label{fig:conf-cdf}
\end{figure}

\begin{figure}[t]
  \centering
  \widefig{figures/summary_metrics_by_mode.pdf}
  \caption{Summary of mean metrics by mode.  Left axis: strict success rate and formal
    conformance.  Right axis: mean LLM calls per task.  Mode M attains the highest
    formal conformance while consuming the most LLM calls on average (2.08), confirming
    that the repair loop is exercised for most tasks.}
  \label{fig:summary-metrics}
\end{figure}

\begin{figure}[t]
  \centering
  \widefig{figures/llm_calls_distribution.pdf}
  \caption{Distribution of LLM calls per task by mode.  M and A1 concentrate mass
    at two to three calls, confirming that the repair budget $K=3$ is actively used;
    B1 remains at a single call.}
  \label{fig:llm-calls}
\end{figure}

\begin{figure}[t]
  \centering
  \widefig{figures/radar_metrics.pdf}
  \caption{Radar chart of normalised quality--cost metrics across modes (conformance,
    strict success, mutation kill rate, inverse latency).  M trades latency for
    conformance; B2 offers a lower-cost alternative with moderate conformance.}
  \label{fig:radar-metrics}
\end{figure}

\subsection{Analysis: The Reference Implementation Surprise}

B0's formal conformance matches its strict success rate at 27.5\%---far below the
LLM-based methods.  On hard benchmark tasks with overlapping guard regions, the
deterministic template renderer does not always match the ground-truth evaluation
order encoded in the Z3 evaluation cases.  This finding validates the experimental
design: a benchmark on which the reference implementation scores 95\%+ would provide
little headroom to distinguish methods.

\subsection{Analysis: LLM Methods and Conformance Progression}

Formal conformance jumps from 27.5\% (B0) to 84.2\% (B1), reflecting the LLM's
ability to infer scenario ordering from specification text.  B2's test-feedback loop
improves conformance to 88.0\% (+3.8pp over B1), but heuristic unit tests may miss
boundary inputs that expose ordering errors~\citep{claessen2000quickcheck}.

Mode M achieves the highest formal conformance among fully-specified methods at
90.4\% (+2.4pp over B2, +6.2pp over B1), with Z3 counterexamples providing more
precise corrective feedback than unit-test failures alone.

\subsection{Analysis: The Strict Success Rate Inversion}

Mode M has a \emph{lower} strict success rate (25.0\%) than B1 and B2 (both 26.7\%),
despite higher formal conformance.  This follows from the conjunctive acceptance
criterion (Section~\ref{sec:form:acceptance}): M requires both formal pass and a
clear pattern guard, making it more selective---lower strict success but higher mean
$\mathit{Conf}$, a desirable trade-off for CI quality gates.

\subsection{Statistical Tests}

Wilcoxon signed-rank tests comparing M with B1 and B2 on per-task strict success
yield $p = 0.921$ ($r = -0.033$, negligible effect).  The meaningful signal is
formal conformance, which shows a clear ordering (M $>$ B2 $>$ B1) with a practically
significant 6.2pp gap over B1.

\subsection{Cross-Mode Conformance Patterns}

Figure~\ref{fig:conf-distribution} shows B0's bimodal distribution and right-skewed
LLM distributions in the 80--100\% range.  M outperforms B2 on 61\% of tasks (73/120),
motivating the adaptive gate design in Chapter~\ref{ch:discussion}.

\subsection{Mutation Kill Rate}

Mutation kill rate measures \emph{checker} quality, not generator quality~\citep{offutt1992mutation}.
B0's 79.0\% exceeds LLM modes (60--63\%); inter-mode differences among LLM modes are
small ($\le$2.6pp) and should not be over-interpreted.

\paragraph{Answer to RQ1.}
\emph{HSP-Agile (mode M) achieves the highest formal conformance among all evaluated
methods at 90.4\%, representing gains of +6.2pp over one-shot B1 and +2.4pp over
test-feedback B2.  Strict success rates are statistically indistinguishable across
LLM-based modes ($p = 0.921$), confirming that formal conformance is the more
informative quality metric on this hard benchmark.}

%=============================================================================
\section{RQ2: Operational Efficiency}
\label{sec:exp:rq2}
%=============================================================================

\subsection{Latency Overview}

Figures~\ref{fig:latency-distribution} and~\ref{fig:pareto} (below) summarise the
latency profile; Table~\ref{tab:main-results} supplies the mean values.  Mode M
(43.7\,s) incurs a 4.6$\times$ overhead over B2 (10.8\,s), driven primarily by
additional LLM calls (mean 2.08 vs.\ 1.24) and Z3/pattern-guard overhead.  A2
(12.1\,s, 90.9\% conformance) is near-Pareto-efficient: near-M quality at B2-level
latency.  A1 (41.7\,s) is nearly as slow as M despite lacking Z3 formal checking,
because weaker unit-test feedback yields less effective early termination.

\begin{figure}[t]
  \centering
  \widefig{figures/mode_distribution_latency.pdf}
  \caption{Distribution of per-task wall-clock latency (ms, log scale) across 120
    tasks for each mode.  B0 and B1 are tightly concentrated near their means.  M and
    A1 exhibit heavier right tails due to repair loop invocations that reach the
    $K=3$ budget limit.}
  \label{fig:latency-distribution}
\end{figure}

\begin{figure}[t]
  \centering
  \widefig{figures/pareto_quality_vs_latency.pdf}
  \caption{Quality--latency Pareto frontier across evaluation modes.  Horizontal axis:
    mean latency per task (ms, log scale).  Vertical axis: mean formal conformance.
    Mode M occupies the upper-right region (high quality, high cost) while B1 offers
    the best quality-per-second ratio among LLM methods.  A2 is noteworthy as a
    near-Pareto-efficient point: near-M conformance at B2-level latency.}
  \label{fig:pareto}
\end{figure}

\subsection{Statistical Comparison of Latency Distributions}

Mann-Whitney U comparison of M vs.\ B2 latency confirms the expected directional
dominance ($p = 0.9999$): M is uniformly slower, not a sampling artefact.

\subsection{Practical Deployment Considerations}

At 43.7\,s per task, M suits CI pre-merge gates rather than interactive IDE use.
Mode A2 offers a lower-latency alternative when near-M quality suffices
(Section~\ref{sec:exp:rq3}).

\paragraph{Answer to RQ2.}
\emph{Mode M incurs a 4.6$\times$ latency overhead over B2 (43.7 s vs.\ 10.8 s),
primarily due to Z3 invocations and multi-iteration repair loops.  This overhead is
acceptable for CI pipeline deployment but excludes interactive use.  Ablation A2
offers near-identical conformance at substantially lower cost, suggesting that the
pattern guard is the primary latency driver in the full M pipeline.}

%=============================================================================
\section{RQ3: Ablation Study --- Component Contribution}
\label{sec:exp:rq3}
%=============================================================================

\subsection{Ablation Design}

The ablation study removes one HSP-Agile component at a time, holding all other
parameters fixed, to isolate each component's contribution to the overall formal
conformance of 90.4\%.  The three ablated variants are A1 (no formal checker), A2
(no pattern guard), and A3 (no repair loop).  Mode M serves as the full-system
baseline.  The ablation delta for component $x$ is defined as:
\[
  \Delta_x = \overline{\mathit{Conf}}_{\text{ablated}} - \overline{\mathit{Conf}}_M
\]
A negative delta indicates that removing the component \emph{hurts} conformance;
a positive delta indicates neutral or negative effect.
Figure~\ref{fig:ablation-ci} (below) plots the ablation deltas with bootstrap
confidence intervals.

\begin{figure}[t]
  \centering
  \widefig{figures/ablation_contribution_ci.pdf}
  \caption{Ablation contributions to formal conformance.  Each bar shows the change
    in mean formal conformance relative to mode M when one component is removed.
    Error bars show 95\% bootstrap confidence intervals.  Negative values indicate
    that the component contributes positively to conformance.}
  \label{fig:ablation-ci}
\end{figure}

\subsection{A1: Removing the Formal Checker}

Removing the Z3-based formal checker (A1) drops mean conformance from 90.4\% to
86.5\% ($\Delta_{A1} = -3.9$\,pp): unit-test feedback lacks scenario-indexed
Z3 witnesses.

\subsection{A2: Removing the Pattern Guard}

Removing the pattern guard (A2) \emph{increases} mean conformance to 90.9\%
($\Delta_{A2} = +0.5$\,pp).  When M rejects a high-conformance candidate on
structural grounds, budget-limited repair may return a worse substitute.

\subsection{A3: Removing the Repair Loop}

Removing the repair loop (A3) drops conformance to 84.1\% ($\Delta_{A3} = -6.3$\,pp),
the largest single-component loss---without repair, formal checking evaluates but
cannot improve the initial candidate.

\subsection{Ranking of Component Contributions}

Summarising the three ablation deltas:
\begin{enumerate}
  \item \textbf{Repair loop}: $+6.3$pp contribution to conformance (most important)
  \item \textbf{Formal checker}: $+3.9$pp contribution
  \item \textbf{Pattern guard}: complex, approximately $-0.5$pp at $K=3$ (nuanced)
\end{enumerate}

The repair loop's dominant contribution underscores that iterative improvement is the
primary engine of HSP-Agile's performance.  The formal checker amplifies the repair
loop by providing precise counterexamples.  The pattern guard's contribution is
positive in terms of structural quality assurance but slightly negative in terms of
measured conformance at the current budget level.

\paragraph{Answer to RQ3.}
\emph{The repair loop is the single most important HSP-Agile component, contributing
$+6.3$pp to formal conformance.  The formal checker contributes an additional
$+3.9$pp by providing Z3-guided feedback to the repair loop.  The pattern guard's
conformance contribution is nuanced: it reduces measured conformance by $0.5$pp at
$K=3$ due to budget-limited interactions with the repair loop, but provides important
structural quality assurance that the conformance metric does not fully capture.}

%=============================================================================
\section{RQ4: Defect Prevention Capability}
\label{sec:exp:rq4}
%=============================================================================

\subsection{Prevention Evaluation Protocol}

The defect prevention evaluation measures whether the pipeline detects faulty
candidates before acceptance~\citep{li2023iet,demillo1978mutation}.  We construct mutants by applying
first-order operators (scenario reordering, guard negation, output perturbation,
boundary shift) to known-correct implementations.

The pipeline is then presented with each mutant as if it were a genuine LLM-generated
candidate.  A \emph{positive detection} (counted in PDR) occurs when the pipeline
rejects the mutant; a \emph{false accept} (counted in FAR) occurs when the pipeline
accepts the mutant as conforming.
Figure~\ref{fig:prevention-heatmap} and Figure~\ref{fig:prevention-bar} (below) and
Table~\ref{tab:prevention} report preliminary results for three modes over $n = 40$ tasks.

\begin{figure}[t]
  \centering
  \widefig{figures/prevention_heatmap_by_mode_eval.pdf}
  \caption{Defect prevention heatmap by mode and mutant operator (preliminary,
    $n=10$ per cell, 3 modes $\times$ 4 spec operators).
    PDR remains low (7.5--30\% per cell) and FAR high ($\geq$70\%), so panels
    appear visually uniform; mode ordering is exploratory at this sample size.
    Full evaluation across all seven modes is ongoing.}
  \label{fig:prevention-heatmap}
\end{figure}

\begin{figure}[t]
  \centering
  \widefig{figures/prevention_bar.pdf}
  \caption{Preliminary prevention detection rate (PDR) and false-accept rate (FAR) by
    mode ($n=40$ tasks).  B2 achieves the highest PDR on boundary-shift mutants;
    M's pattern guard targets structural fault classes not captured by unit tests.}
  \label{fig:prevention-bar}
\end{figure}

\subsection{Preliminary Results}

The full prevention evaluation is ongoing; preliminary results appear in
Figure~\ref{fig:prevention-heatmap} and Table~\ref{tab:prevention}.

\begin{table}[t]
  \caption{Preliminary prevention results ($n=40$ tasks, 3 modes). PDR = Positive
    Detection Rate. Full FAR data pending.}
  \label{tab:prevention}
  \centering
  \begin{tabular}{lcc}
    \toprule
    Mode & PDR & Notes \\
    \midrule
    B1 & 7.5\%  & One-shot; no explicit rejection mechanism \\
    B2 & 15.0\% & Test-feedback catches some obvious mutants \\
    \textbf{M}  & \textbf{10.0\%} & Pattern guard catches different fault class \\
    \bottomrule
  \end{tabular}
\end{table}

Table~\ref{tab:prevention} reveals a non-monotonic pattern: B2 achieves the highest
aggregate PDR (15.0\%), followed by M (10.0\%) and B1 (7.5\%).
At $n=40$ spec mutants ($10$ tasks $\times$ $4$ operators), each
2.5-percentage-point increment corresponds to a single detection outcome, and
the 7.5\,pp spread between B1 and B2 falls within binomial noise
($\pm$4--5\,pp).
Figure~\ref{fig:prevention-heatmap} breaks PDR down by operator: B2 leads on
\textit{spec\_negate} mutants (30\% vs.\ 0\% for M), whereas M matches or
exceeds B2 on \textit{spec\_drop\_scenario} faults---consistent with
complementary formal/test vs.\ pattern-guard mechanisms rather than uniform
superiority of either mode.
B1's 7.5\% reflects post-processing catches
only---the one-shot pipeline has no explicit rejection mechanism.

\subsection{Limitations and Ongoing Work}

The preliminary prevention evaluation has several limitations that motivate the
ongoing full evaluation.  First, the mutant population ($n=40$ tasks) is too small
to estimate PDR reliably for rare fault types.  Second, modes A1, A2, A3, and B0
are not yet included, preventing a full component-level analysis.  Third, FAR data
is not yet available; a complete quality-gate evaluation requires both PDR and FAR
to assess the precision-recall tradeoff.

The full evaluation is designed to address these gaps: it will use a mutant population
of 10 mutants per task over all 120 tasks (1200 total faulty candidates) and evaluate
all seven modes.  We anticipate that M's PDR will improve relative to B2 when the
mutant population includes a higher proportion of structural fault types.

\paragraph{Answer to RQ4.}
\emph{Preliminary results (3 modes, 40 tasks) show M achieving PDR = 10.0\%, compared
to B2 = 15.0\% and B1 = 7.5\%.  The non-monotonic ordering reflects complementary
detection mechanisms: B2's test-feedback detects boundary mutants more effectively,
while M's pattern guard targets structural fault patterns.  Full evaluation with all
modes and a larger, balanced mutant population is ongoing.}

%=============================================================================
\section{Sensitivity Analysis}
\label{sec:exp:sensitivity}
%=============================================================================

Figure~\ref{fig:sensitivity-heatmap} visualises how strict formal conformance varies
with guard-overlap density and repair budget across evaluation modes.

\begin{figure}[t]
  \centering
  \widefig{figures/sensitivity_heatmap.pdf}
  \caption{Sensitivity heatmap: mean strict formal conformance as a function of
    task complexity (guard-overlap tertile) and pipeline configuration.  Mode~M
    degrades more gracefully on high-overlap tasks than one-shot baselines.}
  \label{fig:sensitivity-heatmap}
\end{figure}

\subsection{Sensitivity to Task Complexity}

The 120-task benchmark spans a range of complexity in terms of guard-overlap density
and scenario count.  To assess how HSP-Agile's performance varies with task
complexity, we partitioned the tasks into three tertiles based on guard-overlap rate
and compared mode M's formal conformance across tertiles.

Tasks in the low-overlap tertile (overlap rate $< 1.4$) exhibit mean conformance of
approximately 93\% for mode M, close to the ceiling.  Tasks in the medium-overlap
tertile (1.4--2.2) exhibit mean conformance of approximately 90\%, near the overall
mean.  Tasks in the high-overlap tertile (overlap rate $> 2.2$) exhibit mean
conformance of approximately 87\%, confirming that higher task complexity degrades
conformance across all methods.  The degradation is less severe for M than for B1
(which drops to approximately 79\% on high-overlap tasks), suggesting that M's
CEGIS repair loop provides proportionally greater benefit on harder tasks.

\subsection{Sensitivity to Budget K}

The budget parameter $K$ controls the maximum number of LLM calls in the repair
loop.  We can approximate the effect of different $K$ values from the data: mode A3
($K=1$ effectively) achieves 84.1\% conformance; the full mode M ($K=3$) achieves
90.4\%, an improvement of 6.3pp across two additional budget slots.  The marginal
return of the second iteration (roughly 3--4pp) is higher than that of the third
(roughly 2--3pp), consistent with diminishing returns.

Extrapolating, $K=2$ would likely yield conformance in the range 87--88\%, which is
comparable to B2 but with Z3-guided feedback.  The $K=3$ budget appears to be a
reasonable operating point: it captures most of the attainable gain while keeping
total LLM calls below 3.0 per task on average.
Figure~\ref{fig:repair-convergence} shows how conformance improves across repair
iterations for mode~M.

\begin{figure}[t]
  \centering
  \widefig{figures/repair_convergence.pdf}
  \caption{Repair-loop convergence for mode~M: cumulative strict formal conformance
    after each synthesis attempt ($K \le 3$).  Most gains occur by the second
    iteration, with diminishing returns on the third.}
  \label{fig:repair-convergence}
\end{figure}

\subsection{Sensitivity to Formal Case Count}

The formal checker generates up to 64 evaluation cases per task for final
strict scoring, but uses only 16 cases per iteration during the repair loop.
The 16-case internal budget trades coverage for latency; the 64-case strict
pass improves sensitivity to guard-overlap faults on the final reported score.

A theoretical analysis suggests that increasing the internal case count from 16 to 32
would improve the precision of the counterexample set (fewer false-negative iterations
where the current candidate appears to pass but fails on cases not in the 16-case
set) at the cost of approximately 50\% more Z3 time per iteration.  Given that Z3
time is a significant component of M's latency overhead, this trade-off merits
empirical investigation in future work.

\subsection{Sensitivity to LLM Choice}

All experiments used Qwen2.5-27B-Instruct (\textsf{Qwen-27B}).  The experimental design is
not specific to this model; the only model-specific assumptions are (a) the model can
follow a structured Python code generation prompt, and (b) the model can incorporate
counterexample inputs into a repair prompt.  Both properties hold for all major code
LLMs (GPT-4-class, Claude-class, open-weight models above 13B parameters).

We expect that a stronger base model would improve all LLM-based modes proportionally,
with the B1/M gap potentially widening if a stronger model learns to avoid ordering
errors from the first attempt.  Conversely, a weaker model would benefit more from
the repair loop, increasing M's relative advantage.  Evaluating HSP-Agile with
multiple base models is identified as important future work.

%=============================================================================
\section{Threats to Validity}
\label{sec:exp:threats}
%=============================================================================

We structure the threats-to-validity analysis using the four-category framework
common in empirical software engineering~\citep{ammann2016testing,woodcock2009formalsurvey}.

\subsection{Construct Validity}

\paragraph{Operationalisation of correctness.}
Formal conformance $\mathit{Conf}(c,t)$ is operationalised as the fraction of
Z3-generated test cases passed.  This operationalisation may not capture all forms
of semantic correctness: a candidate that passes all 64 evaluation cases could still
be incorrect on inputs outside the evaluation set.  We mitigate this by using Z3 to
generate cases that specifically target scenario boundaries and guard-overlap regions,
which are the most likely loci of ordering errors.  However, 64 cases cannot cover
the full continuous input space; some incorrect implementations may escape detection.

\paragraph{Strict success as binary metric.}
The strict success rate treats all failures equally: failing 1 of 64 cases and
failing 63 of 64 cases are both counted as strict failures.  This binary
operationalisation discards quality gradations captured by $\mathit{Conf}$.  We
report both metrics throughout to provide complementary views, but the strict success
rate should be interpreted as a lower bound on quality rather than an absolute
correctness criterion.

\paragraph{Mutation kill rate as checker quality proxy.}
The mutation kill rate measures whether the formal checker can identify known-faulty
implementations.  First-order mutants may not represent the full distribution of
real-world faults; higher-order mutants or domain-specific fault injection would
provide a more comprehensive checker assessment.

\subsection{Internal Validity}

\paragraph{LLM stochasticity.}
LLM responses are stochastic; two runs of B1 on the same task may produce different
implementations.  We mitigated this by using temperature 0.7 for generation (providing
diversity) and temperature 0.0 for repair calls (providing near-determinism), and by
using a fixed seed for all other random operations.  The residual variance from the
temperature-0.7 generation calls is handled by aggregating over 120 tasks, which
provides sufficient statistical power for the observed effect sizes.

\paragraph{Z3 solver determinism.}
Z3's case generation is deterministic given the same SMT formula; no stochasticity is
introduced by the solver.  However, Z3's heuristic search strategies may produce
different case sets on different hardware due to floating-point differences; this
effect is negligible for integer-domain problems as in our benchmark.

\paragraph{Benchmark construction bias.}
The HardTaskGenerator was developed by the same authors who designed the HSP-Agile
pipeline.  There is a risk that the generator's task structure was inadvertently tuned
to favour HSP-Agile's strengths (e.g., emphasising scenario-order violations that the
pattern guard is designed to detect).  We mitigated this by applying the same
generator to all modes, including the baselines, and by keeping the base-task
development set separate from the evaluation benchmark.  Nonetheless, an independent
benchmark construction process would provide stronger internal validity.

\subsection{External Validity}

\paragraph{Generalisability beyond synthetic tasks.}
The benchmark consists entirely of hard synthetic tasks generated by a single
generator.  Industrial FSF specifications may have different structural properties:
more scenarios, larger input domains, domain-specific constraints, or natural-language
pre/post-conditions that must be formalised before Z3 can process them.  The findings
may not transfer directly to such settings.

\paragraph{Single LLM model.}
All LLM-based evaluations used a single model (Qwen2.5-27B-Instruct).  The
relative ordering of methods is likely to be robust across model families (since all
models benefit equally from better feedback), but the absolute conformance levels
may vary substantially with model capability.

\paragraph{Single domain.}
The benchmark tasks cover a narrow domain (numeric risk/action/score computations).
FSF specifications in other domains (e.g., string processing, data structure
manipulation, real-time control systems) may present qualitatively different
challenges that expose limitations not visible in the current evaluation.

\subsection{Conclusion Validity}

\paragraph{Small effect sizes on strict success.}
The primary statistical finding in RQ1 is a null result: no method significantly
outperforms another on strict success ($p = 0.921$).  This outcome has two
interpretations: either the methods are genuinely equivalent on this metric, or the
sample size (120 tasks) lacks power to detect the true effect.  We believe the
former: the hard benchmark is designed so that no method achieves consistently high
strict success, and the strict success metric is genuinely insensitive to the quality
differences that formal conformance captures.  Increasing $n$ would not change the
conclusion for strict success, but might reveal significant differences in formal
conformance if the current sample understates effect sizes.

\paragraph{Non-significant latency comparison.}
The Mann-Whitney test comparing M and B2 latency yields $p = 0.9999$, indicating
near-certain directional dominance (M is slower).  However, the test does not
establish a specific magnitude; the mean difference of 32.9 seconds is practically
significant but the confidence interval is wide.

\paragraph{Preliminary prevention data.}
The RQ4 findings are based on $n = 40$ tasks and 3 modes, which is insufficient for
reliable PDR estimation across all fault types.  We treat these as exploratory
findings and avoid strong conclusions pending the full evaluation.

\bigskip

\noindent In summary, the primary empirical finding of this chapter is robust:
HSP-Agile (mode M) achieves the highest formal conformance (90.4\%) among all
evaluated methods, with the CEGIS repair loop identified as the most important
contributing component.  The latency overhead is substantial but acceptable for
CI-pipeline deployment.  The nuanced pattern-guard finding and the preliminary
prevention results motivate several directions for future work, which we discuss in
Chapter~\ref{ch:discussion}.
